\documentclass[12pt,a4paper]{article}
\usepackage[utf8]{inputenc}
\usepackage[T1]{fontenc}
\usepackage[brazil]{babel}
\usepackage[left=3cm,right=2cm,top=3cm,bottom=2cm]{geometry}
\usepackage{fancyhdr}
\usepackage{setspace}
\usepackage{parskip}
\usepackage{array}
\usepackage{booktabs}
\usepackage{enumitem}

% ==========================================
% CABEÇALHO E RODAPÉ
% ==========================================
\pagestyle{fancy}
\fancyhf{}
\fancyhead[C]{%
  \textbf{\large AUTORIDADE PORTUÁRIA}\\
  \small Diretoria Técnica e Operacional
}
\fancyfoot[C]{\small Página \thepage}
\fancyfoot[L]{\small Documento gerado pelo PortoGpt}
\renewcommand{\headrulewidth}{0.4pt}
\renewcommand{\footrulewidth}{0.4pt}

% ==========================================
% INÍCIO DO DOCUMENTO
% ==========================================
\begin{document}
\onehalfspacing

\begin{center}
  \vspace*{1cm}
  {\LARGE \textbf{RELATÓRIO TÉCNICO}}\\[0.5cm]
  {\large Nº 001/2026 -- DTO}\\[1cm]
\end{center}

% --- Dados do Relatório ---
\begin{tabular}{@{}ll@{}}
  \textbf{Data:} & 07 de maio de 2026 \\
  \textbf{Setor:} & Diretoria Técnica e Operacional \\
  \textbf{Responsável:} & Eng. Maria Oliveira \\
  \textbf{Assunto:} & Inspeção das Condições Estruturais do Cais 03 \\
\end{tabular}

\vspace{1cm}
\hrule
\vspace{0.5cm}

\section{Objetivo}

O presente relatório tem por objetivo apresentar os resultados da inspeção técnica realizada nas instalações do Cais 03 desta Autoridade Portuária, com foco na avaliação das condições estruturais e operacionais.

\section{Metodologia}

A inspeção foi realizada no dia 05 de maio de 2026, por equipe técnica composta por engenheiros civis e inspetores de segurança. Foram utilizados os seguintes métodos:

\begin{itemize}[leftmargin=2cm]
  \item Inspeção visual detalhada;
  \item Ensaios de esclerometria nas estruturas de concreto;
  \item Medição de potencial de corrosão nas armaduras;
  \item Levantamento topográfico dos recalques.
\end{itemize}

\section{Resultados}

\subsection{Condições Estruturais}

As estruturas de concreto armado do cais apresentam, de modo geral, boas condições de conservação. Foram identificados os seguintes pontos de atenção:

\begin{table}[h]
\centering
\begin{tabular}{@{}clc@{}}
  \toprule
  \textbf{Item} & \textbf{Descrição} & \textbf{Severidade} \\
  \midrule
  1 & Fissuras superficiais no trecho norte & Baixa \\
  2 & Desgaste do pavimento na área de manobra & Média \\
  3 & Corrosão leve nos defensores metálicos & Média \\
  4 & Recalque de 2mm na junta de dilatação J3 & Baixa \\
  \bottomrule
\end{tabular}
\caption{Achados da inspeção técnica}
\end{table}

\section{Conclusões e Recomendações}

Com base nos resultados obtidos, recomenda-se:

\begin{enumerate}[leftmargin=2cm]
  \item Programar reparo das fissuras no trecho norte no próximo período de estiagem;
  \item Substituir o pavimento degradado na área de manobra;
  \item Aplicar tratamento anticorrosivo nos defensores metálicos;
  \item Monitorar trimestralmente o recalque na junta J3.
\end{enumerate}

\vspace{1.5cm}

\begin{center}
\textbf{Eng. Maria Oliveira}\\
CREA-MA 12345/D\\
Diretoria Técnica e Operacional
\end{center}

\end{document}
